On simulated random-neighbour NK landscapes at two $(L,A)$ settings, additive ridge absorbs $K=0$ completely, and pairwise ridge absorbs $K=1$ at $(15,4)$ ($0.998$ Spearman at $N=1000$) but not at $(20,20)$ ($0.138$); since only these two settings were run, the difference cannot be attributed to alphabet size rather than length. At $K=4$, and at $(20,20)$ with $K\ge2$, no model learned much from $N=1000$ (mean Spearman at most $0.080$); larger $N$ was not tested in those regimes. The small CNN and MLP did not beat pairwise ridge at $N=1000$ and stayed below it up to $N=4000$ on the two sweep tasks, with a narrowing gap on L15\_A4\_K2 only. The registered design criterion was met in sign, but the association between prediction gains and single-step design gains rests on one task cell, and with 10 proposals and five seeds the design hit rates at $K\ge2$ could not be reliably separated from chance. A fuller study would add more seeds and larger proposal batches for the design endpoint, intermediate $K$, $L$ and $A$ varied separately, noisy labels, correlated (non-i.i.d.) landscapes, stronger neural baselines, larger $N$ in the hard regimes, and real assays; none of that was done here.
